% =====================================================================
%  RESUMO + ABSTRACT
% =====================================================================
\chapter*{Resumo}
\addcontentsline{toc}{chapter}{Resumo}
\markboth{Resumo}{Resumo}

O arquivo familiar — fotografias, documentos e árvores genealógicas
acumulados ao longo de gerações — encerra um enorme valor afetivo e
histórico, mas permanece tipicamente disperso, desorganizado e, acima de
tudo, \emph{não narrado}. Este projeto propõe e implementa um sistema que
transforma automaticamente esse material num produto narrativo
multimédia, sem exigir que o utilizador escreva textos ou descreva o
conteúdo manualmente.

O sistema organiza-se num \emph{pipeline} de quatro módulos sequenciais:
\textbf{(M1)} ingestão multimodal, que extrai metadados \textsc{exif},
reconhece texto em documentos por \textsc{ocr} e descreve o conteúdo
visual das fotografias com um modelo de visão; \textbf{(M2)} organização
temporal, que resolve datas incompletas e constrói um grafo de
parentesco; \textbf{(M3)} geração narrativa, assente em modelos de
linguagem de grande escala (\textsc{llm}) orientados por
\emph{Retrieval-Augmented Generation} (\textsc{rag}) e por um conjunto de
\emph{templates} calibrados para português europeu; e \textbf{(M4)}
geração multimédia, que converte a narrativa num vídeo documental com
enquadramento cinematográfico, narração por síntese de voz e legendas
sincronizadas. A narrativa é segmentada em \emph{cenas} ligadas a
fotografias concretas, o que permite ao M4 sincronizar cada imagem com a
sua narração.

A solução foi desenvolvida como uma aplicação Web completa, com um
\emph{backend} assíncrono em \textsc{FastAPI} e um \emph{frontend} em
\textsc{React}/\textsc{TypeScript}, e evoluiu de um protótipo
\emph{local-first} para uma arquitetura cliente-servidor multi-inquilino
implantada na nuvem (\textsc{Render}, \textsc{Vercel} e \textsc{Supabase}
para autenticação, base de dados e armazenamento), correndo de forma
idêntica em ambiente local e \emph{online}. Adota uma estratégia de
inteligência artificial \textbf{em cascata e a custo nulo} — texto via
\emph{Groq} (Llama~3.3~70B, na nuvem), com a inferência local via
\emph{Ollama} e o \emph{Google~Gemini} como alternativas, e visão via
\emph{Gemini} — e técnicas de engenharia de robustez (degradação graciosa
do \textsc{rag}, \emph{fallbacks} de modelo e rebaixamento síncrono) que
mantêm o sistema operacional mesmo em ambientes com dependências
reduzidas. Os resultados
demonstram a viabilidade de gerar, de forma autónoma e a custo
essencialmente nulo, narrativas factualmente ancoradas e vídeos
documentais coerentes a partir de um arquivo familiar heterogéneo.

\vspace{0.8cm}
\noindent\textbf{Palavras-chave}\\[0.2cm]
Inteligência Artificial Generativa; Modelos de Linguagem de Grande
Escala; Retrieval-Augmented Generation; Processamento Multimodal;
Genealogia Computacional; Geração Automática de Vídeo; Narrativa Digital;
FastAPI; React.

\chapter*{Abstract}
\addcontentsline{toc}{chapter}{Abstract}
\markboth{Abstract}{Abstract}

The family archive — photographs, documents and genealogical trees
accumulated across generations — holds enormous emotional and historical
value, yet it typically remains scattered, disorganised and, above all,
\emph{untold}. This project proposes and implements a system that
automatically turns such material into a multimedia narrative product,
without requiring the user to write any text or describe the content
manually.

The system is organised as a pipeline of four sequential modules:
\textbf{(M1)} multimodal ingestion, extracting \textsc{exif} metadata,
running \textsc{ocr} on documents and describing the visual content of
photographs with a vision model; \textbf{(M2)} temporal organisation,
resolving incomplete dates and building a kinship graph; \textbf{(M3)}
narrative generation, based on Large Language Models steered by
Retrieval-Augmented Generation and a set of European-Portuguese prompt
templates; and \textbf{(M4)} multimedia generation, turning the narrative
into a documentary video with cinematic framing, text-to-speech
narration and synchronised captions.

The solution was built as a complete web application, with an
asynchronous \textsc{FastAPI} backend and a \textsc{React}/\textsc{TypeScript}
frontend, evolving from a local-first prototype into a cloud-deployed,
multi-tenant client–server architecture (Render, Vercel and Supabase for
authentication, database and storage), running identically in local and
online environments. It adopts a \textbf{zero-cost, cascading AI strategy}
— text via Groq (cloud Llama~3.3~70B), with local Ollama inference and
Google~Gemini as fallbacks, and vision via Gemini — and robustness
engineering techniques (graceful RAG degradation, model fallbacks and
synchronous downgrade) that keep the system operational even in
dependency-constrained environments. The
results demonstrate the feasibility of autonomously generating, at
essentially zero cost, factually grounded narratives and coherent
documentary videos from a heterogeneous family archive.

\vspace{0.8cm}
\noindent\textbf{Keywords}\\[0.2cm]
Generative Artificial Intelligence; Large Language Models;
Retrieval-Augmented Generation; Multimodal Processing; Computational
Genealogy; Automatic Video Generation; Digital Storytelling; FastAPI;
React.
